\subsection{Первое задание.}
\begin{problem}[4]
    Два игрока играют в честную игру (с вероятностью победы в каждой из партий 1/2).
    Игроки договорились, что тот, кто первым выиграет 6 партий, получит весь приз.
    Предположим, что игроки не успели доиграть: первый выиграл 5 партий, а второй -- 3.
    Как следует разделить приз?
\end{problem}
\begin{solution}
    Приз следует разделить в соответствующей пропорции с шансами выигрыша: поскольку второй игрок выиграет только если он победит три раза подряд, то есть с вероятностью, равной 1/8, а с оставшимися 7/8 победит первый игрок, то приз нужно поделить в отношении 7:1.
\end{solution}
\begin{problem}[7]
    Из урны, содержащей $a$ белых и $b$ чёрных шаров, извлекается наугад один шар и откладывается в сторону.
    Какова вероятность того, что извлечённый наугад второй шар окажется белым, если:
    \begin{enumerate}
        \item первый извлечённый шар белый;
        \item цвет первого извлечённого шара остаётся неизвестным?
    \end{enumerate}
\end{problem}
\begin{solution} \leavevmode
    \begin{enumerate}
        \item Распишем по определению и проведём несложные арифметические преобразования:
        \[
            \P(2 - white \mid 1 - white) = \dfrac{\P(2 - white \wedge 1 -  white)}{\P(1 - white)} = \dfrac{\dfrac{a}{a + b}\cdot \dfrac{a - 1}{a + b - 1}}{\dfrac{a}{a + b}} = \dfrac{a - 1}{a + b - 1}.
        \]
        \item Используем формулу полной вероятности и получим:
        \begin{multline*}
            \P(2 - white) = \P(2 - white \mid 1 - white)\P(1 - white) + \P(2 - white \mid 1 - black)\P(1 - black) = \\ =
            \dfrac{(a - 1)a}{(a + b - 1)(a + b)} + \dfrac{\dfrac{b}{a + b}\cdot \dfrac{a}{a + b - 1}}{\dfrac{b}{a + b}}\dfrac{b}{a + b} = \dfrac{a(b + a - 1)}{(a + b)(a + b - 1)} = \dfrac{a}{a + b}.
        \end{multline*}
    \end{enumerate}
\end{solution}
\begin{problem}[17]
    В гардеробе случайным образом перепутались $N$ одинаковых шляп посетителей.
    Какова вероятность того, что хотя бы один посетитель получит свою шляпу?
    Рассмотреть случаи $N = 4, N = 10000$.
\end{problem}
\begin{solution}
    Заметим, что события $A:$ хотя бы посетитель получил свою шляпу и $B:$ никто не получил свою шляпу несовместны, и $\P(A) + \P(B) = 1$.
    Поскольку $B = \bigcup\limits_{i = 1}^n B_i$, где $B_i$~---~$i$-й посетитель не получил свою шляпу, то согласно формуле включений исключений
    \begin{multline*}
        \P(A) = 1 - \P(B) = 1 - \P(\bigcup\limits_{i = 1}^n B_i) = \\ = 1 - \dfrac{1}{n!}(n! - \binom{n}{1}(n - 1)! + \ldots + (-1)^n\binom{n}{n}) = 1 - \sum\limits_{i = 0}^n(-1)^j\binom{n}{j}\dfrac{(n - j)!}{n!} = \\ = 1 - \sum\limits_{j = 0}^n \dfrac{(-1)^j}{j!} \ra 1 - e^{-1}, n \ra +\infty.
    \end{multline*}
    Для $N = 4$ вероятность составляет $5/8$.
\end{solution}
\begin{problem}[30]
    Опыт состоит в подбрасывании симметричной монеты до тех пор, пока два раза подряд она не выпадет одной и той же стороной.
    Подбрасывания независимы в совокупности.
    Постройте пространство элементарных событий и найдите вероятности следующих событий:
    \begin{enumerate}
        \item опыт закончится до шестого бросания;
        \item для завершения опыта потребуется чётное число бросаний.
    \end{enumerate}
\end{problem}
\begin{solution}
    Пусть $\Omega$~---~множество финитных бинарных последовательностей с двумя единицами на конце.
    Если $A^1_n$ и $A^2_n$~---~это такие последовательности бросков длины $n$, одна из которых начинается с 0, а вторая с 1, то присвоим им вероятность $\P(\{A^1_n\}) = \P(\{A^2_n\}) = 1/2^n$.
    Таким образом определено вероятностное пространство на $(\Omega, 2^{\Omega}, \P)$ (где мы продолжаем по корректности последовательности бросков).
    Пусть $A_i = \{\text{опыт закончится на }i\text{-ом шаге}\}$.
    Тогда $\P(A_i) = \dfrac{2}{2^i} = \dfrac{1}{2^{i - 1}}, i \geq 2$.
    \begin{enumerate}
        \item Первый пункт -- просто $\P\biggr(\bigcup\limits_{i = 2}^{6} A_i\biggr) = 1/2 + 1/4 + \ldots 1/16 = 1 - 1/16 = 15/16$.
        \item В этом случае у нас искомая вероятность равна
        \[
            \P\biggr(\bigcup_{i = 1}^{+\infty}A_{2i}\biggr) = \sum\limits_{i = 1}^{+\infty} \dfrac{1}{2^{2i - 1}} = 2\sum\limits_{i = 1}^{+\infty}\dfrac{1}{4^i} = \dfrac{2}{3}.
        \]
    \end{enumerate}
\end{solution}
\begin{problem}[37]
    Покажите, что из независимости событий $A$ и $B$ следует независимость событий $A$ и $\overline{B}$, $\overline{A}$ и $B$, $\overline{A}$ и $\overline{B}$.
\end{problem}
\begin{solution}
    Пусть $A$ и $B$~---~независимые события. Тогда, $\P(\overline{A}B) = \P(B) - \P(AB) = \P(B) - \P(A)\P(B) = \P(B)(1 - \P(A)) = \P(B)\P(\overline{A})$ что и показывает независимость $\overline{A}$ и $B$.
    В силу коммутативности $\R$ по умножению, если $A$ не зависит от $B$, то и $B$ не зависит от $A$, а следовательно утверждение о независимости $\overline{B}$ и $A$ является частным случаем предыдущего случая. А теперь применим первый случай к $B$ и $\overline{A}$ и получим, что из независимости $B$ и $\overline{A}$ следует независимость $\overline{B}$ и $\overline{A}$, а поскольку из независимости $A$ и $B$ следует независимость $B$ и $\overline{A}$, то и третья импликация показана.
\end{solution}
\begin{problem}[47]
    Известно, что $96\%$ выпускаемой продукции соответствует стандарту.
    Упрощённая схема контроля признаёт годным с вероятностью $0.98$ каждый стандартный экзмепляр аппаратуры и с вероятностью $0.05$ -- каждый нестандартный экземпляр аппаратуры.
    Найдите вероятность того, что изделие, прошедшее контроль, соответствует стандарту.
\end{problem}
\begin{solution}
    Пусть $A$~---~событие, что изделие прошло контроль, а $B$~---~что изделие соответствует стандарту.
    Необходимо найти $\P(B \mid A)$.
    Используем формулы Байеса и полной вероятности и получим:
    \[
        \P(B \mid A) = \dfrac{\P(A \mid B)\P(B)}{\P(A \mid B)\P(B) + \P(A \mid \overline{B})\P(\overline{B})} = \dfrac{0.96 \cdot 0.98}{0.96 \cdot 0.98 + 0.05 \cdot 0.04} \approx 0.998.
    \]
\end{solution}
\begin{problem}[50]
    Пусть случайная величина $X$ имеет показательное распределение $Exp(\lambda)$.
    Покажите, что имеет место <<отсутствие последействия>>:
    \[
        \P(X > x + y \mid X > x) = \P(X > y).
    \]
    Верно ли обратное утверждение?
    Приведите пример дискретной случайной величины со свойством <<отсутствия последействия>>.
\end{problem}
\begin{solution}
    Если $X \sim Exp(\lambda)$, то $\P(X \leq x) = 1 - e^{-\lambda x}. x \geq 0$ (если $x < 0$, то равно 0).
    Тогда
    \begin{multline*}
        \P(X > x + y \mid X > x) = \dfrac{1 - \P(X \leq x + y)}{1 - \P(X \leq x)} = \\ =
        \dfrac{1 - (1 - e^{-\lambda (x + y)})}{1 - (1 - e^{-\lambda x})} = e^{-\lambda y} =  \\ =
        1 - (1 - e^{-\lambda y}) = \P(X > y).
    \end{multline*}
    Неверно, что всякая величина, имеющая свойство <<отсутствия последействия>> распределена экспоненциально. \\
    Пусть $Y \sim Geom(p)$, тогда её функцией распределения является $F_Y(n) = 1 - (1 - p)^n$ (при функции вероятности $\P(Y = k) = (1 - p)^{k - 1}p, k \in \N$).
    Покажем, что $Y$ также обладает свойством отсутствия последействия:
    \begin{multline*}
        \P(Y > x + y \mid Y > x) = \dfrac{1 - \P(Y \leq x + y)}{1 - \P(Y \leq x)} = \\ =
        \dfrac{1 - (1 - (1 - p)^{x + y})}{1 - (1 - (1 - p)^{x})} = (1 - p)^{y} = 1 - (1 - (1 - p)^y) =  \\ =
        1 - \P(Y \leq y) = \P(Y > y).
    \end{multline*}
\end{solution}
\begin{problem}[54]
    Пусть $X$ и $Y$~---~независимые случайные величины, имеющие распределение Пуассона.
    Доказать, что случайная величина $Z = X + Y$ имеет распределение Пуассона.
    Найдите вид условного распределения случайной величины $X$ при фиксированном значении случайной величины $Z$.
\end{problem}
\begin{solution}
    Пусть $X \sim Poiss(\lambda_1), Y \sim Poiss(\lambda_2)$~---~независимые случайные величины.
    Тогда, поскольку $\phi_{Poiss(\lambda)}(t) = \exp(\lambda(e^{it} - 1))$, то для $X + Y$ выполняется что $\phi_{X + Y} = \phi_X \cdot \phi_Y = \exp((\lambda_1 + \lambda_2)(e^{it} - 1))$ и значит $X + Y \sim Poiss(\lambda_1 + \lambda_2)$. \\
    По определению условного распределения,
    \begin{multline*}
        \P(X = k \mid Z = n) = \dfrac{\P(x = k \wedge Z = n)}{\P(Z = n)} = \dfrac{\P(X = k \wedge Y = n - k)}{\P(Z = n)} = \\
        = \dfrac{\P(X = k)\P(Y = n - k)}{\P(Z = n)} = \dfrac{\dfrac{e^{-\lambda_1 - \lambda_2}\lambda_1^k\lambda_2^{n - k}}{k!(n - k)!}}{\dfrac{e^{-\lambda_1 - \lambda_2}(\lambda_1 + \lambda_2)^n}{n!}} = \binom{n}{k}\dfrac{\lambda_1^k\lambda_2^{n - k}}{(\lambda_1 + \lambda_2)^n}.
    \end{multline*}
    и условное распределение $\sim Binom(n, \dfrac{\lambda_1}{\lambda_1 + \lambda_2})$.
\end{solution}
\begin{problem}[75]
    Имеется $n$ пронумерованных писем и $n$ пронумерованных конвертов.
    Письма случайным образом раскладываются по конвертам (все $n!$ способов равновероятны).
    Найдите математическое ожидание числа совпадения номеров письма и конвертов (письмо лежит в конверте с тем же номером).
\end{problem}
\begin{solution}
    Пусть $1, 2, \ldots, n$~---~письма и $1, 2, \ldots, n$~---~конверты.
    Тогда, если мы обозначим за $\xi_i = I_{\{i\text{-ое письмо попало в }i\text{-ый конверт.}\}}$, то обозначая за $X = \sum\limits_{i = 1}^n \xi_i$~---~число совпадений номеров письма и конверта, мы получаем
    \[
        \Ex X = \Ex \sum\limits_{ i = 1}^n \xi_i = \sum\limits_{i = 1}^n \Ex \xi_i = \sum\limits_{i = 1}^n \dfrac{1}{n} = 1.
    \]
\end{solution}
\begin{problem}[80]
    Покажите, что если независимые случайные величины $X_1, \ldots, X_n$ имеют показательное распределение с $\lambda_1, \ldots, \lambda_n$ соответственно, то $\min\{X_1, \ldots, X_n\}$ также распределено показательно, причём с параметром $\sum\limits_{k = 1}^n \lambda_k$.
\end{problem}
\begin{solution}
    Уже решено ранее, а если быть точнее то это первая задача пятой недели.
\end{solution}
\begin{problem}[82]
    Допустим, что вероятность столкновения молекулы с другими молекулами в промежутке времени $[t, t + \Delta t)$ равна $p = \lambda \Delta t + o(\Delta t)$ и не зависит от времени, прошедшего после предыдущего столкновения $(\lambda = const)$.
    Найдите распределение времени свободного пробега молекулы и вероятность того, что это время превысит заданную величину $t^*$.
\end{problem}
\begin{solution}
    Разобъём полуинтервал $\Delta = [0, t)$ на $n$ равных полуинтервалов $\Delta_i = [\dfrac{it}{n}, \dfrac{(i + 1)t}{n}), i = \overline{0, n - 1}$.
    Тогда если $A_i$~---~событие о столкновении на $i$-ом полуинтервале, то $\P(A_i) = \lambda t/n + o(t/n)$~---~по условию.
    Отсутствие столкновений на всём полуинтервале $B$ равносильно отсутствию столкновений на каждом полуинтервале, то есть
    \[
        \P(B) = \P(\bigcap\limits_{i = 0}^{n - 1}\overline{A_i}) = (1 - (\lambda t/n + o(t/n)))^n \ra e^{-\lambda t}, n \ra +\infty.
    \]
    Таким образом, время свободного пробега распределено экспоненциально и $\P(t > t^*) = e^{-\lambda t^*}$.
\end{solution}
\begin{problem}[99]
    Спортсмен стреляет по круговой мишени.
    Вертикальная и горизонтальная координаты точки попадания пули (при условии, что центр мишени -- начало координат) -- независимые случайные величины, каждая с распределением $N(0, 1)$.
    Покажите, что расстояние от точки попадания до центра имеет плотность распределения вероятностей $r \exp(-r^2/2)$ для $r \geq 0$.
    Найдите медиану этого распределения.
\end{problem}
\begin{solution}
    Заметим, что случайный вектор $(X, Y)^T$~---~гауссовский (поскольку $X$ и $Y$~---~независимые) и имеет распределение $(X, Y)^T \sim N(\vec{0}, I)$.
    Плотность совместного распределения имеет вид
    \[
        p_{X, Y}(x, y) = \dfrac{1}{2\pi}\exp\biggr(-\dfrac{x^2 + y^2}{2}\biggr).
    \]
    Рассмотрим отображение $g: (r, \phi) \mapsto (x, y)$, где $x = r\cos\phi$ и $y = r\sin\phi$.
    Его дифференциал
    \[
        Dg =
        \begin{pmatrix*}
            \cos\phi & -r\sin\phi \\ \sin\phi & r\cos\phi
        \end{pmatrix*}
    \]
    Тогда плотность в полярных координатах будет выражаться как
    \[
        p_{R, \Phi}(r, \phi) = \dfrac{1}{\det Dg^{-1}}f_{X, Y}(g(r, \phi)) = |\det Dg|f_{X, Y}(r\cos\phi, r\sin\phi) = \dfrac{r}{2\pi}\exp\biggr(-\dfrac{r^2}{2}\biggr), r \geq 0.
    \]
    Для получения плотности $\R$ необходимо проинтегрировать плотность по $\phi$:
    \[
        p_{R}(r) = \int_{-\pi}^\pi p_{R, \Phi}(r, \phi)d\phi = r\exp\biggr(-\dfrac{r^2}{2}\biggr), r \geq 0.
    \]
    Найдём медиану этого распределения. Условием медианы $m$ является
    \[
        \int_{-\infty}^m p_R(r)dr = 1/2.
    \]
    То есть
    \[
        \int_0^m re^{-r^2/2}dr = -\int_0^m e^{-r^2/2}d(-r^2/2) = -(e^{-m^2/2} - 1) = 1 - e^{-m^2/2}.
    \]
    Тогда $1/2 = 1 - e^{-m^2/2}$ и мы получаем, что $m = \sqrt{2\ln 2}$.
\end{solution}